\section{Presupuesto}

Presupuesto estimado en pesos colombianos a precios de 2026. Los materiales los financia el estudiante; el costo del personal se valora pero no representa desembolso. La Universidad aporta en especie taller, laboratorios e instrumentación.

\begin{table}[H]
\centering
\singlespacing
\scriptsize
\caption{Presupuesto estimado del proyecto (COP, 2026).}
\label{tab:presupuesto}
\renewcommand{\arraystretch}{0.95}
\begin{tabular}{p{8.6cm} c r}
\toprule
\textbf{Rubro} & \textbf{Fuente} & \textbf{Valor (COP)}\\
\midrule
Estudiante: \num{384} h (\num{16} semanas $\times$ \num{24} h) a \cop{25000}/h & Especie & \cop{9600000}\\
Asesor: \num{32} h a \cop{90000}/h & EAFIT & \cop{2880000}\\
\midrule
Subsistema térmico: resistencia tubular, canal reflector, aislamiento, termopares tipo K y MAX31855 & Estudiante & \cop{280000}\\
Control e interfaz: Arduino Mega 2560, SSR \SI{25}{\ampere} y pantalla táctil HMI de 4,3 pulgadas & Estudiante & \cop{430000}\\
Accionamiento: motor NEMA 23, reductor, \emph{driver} y encoder & Estudiante & \cop{450000}\\
Estructura, transmisión y mecanizado: perfilería, lámina, tornillería, rodamientos, ejes y acoples & Estudiante & \cop{950000}\\
Fuente \SI{24}{\volt}, gabinete, cableado, protecciones y paro de emergencia & Estudiante & \cop{300000}\\
Acrílico para ensayos (\num{3}, \num{4.5} y \SI{6}{\milli\meter}) e instrumentación de validación & Estudiante & \cop{450000}\\
Imprevistos (\SI{10}{\percent} de materiales) & Estudiante & \cop{286000}\\
\midrule
\textbf{Subtotal personal} & & \textbf{\cop{12480000}}\\
\textbf{Subtotal materiales e insumos} & & \textbf{\cop{3146000}}\\
\textbf{Total} & & \textbf{\cop{15626000}}\\
\bottomrule
\end{tabular}
\end{table}
